\chapter{Conclusion and Future Scope}
\label{ch:conclusion}

\section{Conclusion}

\todo{write last. State what was built, which objectives from
Chapter~\ref{ch:introduction} were met, and what the demonstration showed ---
without overclaiming. In particular, do not claim optimality}

\section{Contributions}

\todo{three or four specific, defensible points}

\section{Limitations}

\begin{itemize}[nosep]
  \item The plan is feasible and explainable, but not globally optimal.
  \item The scenario is fixed and small: one factory, two products, two
        orders.
  \item The database and primary language model are cloud services, so the
        system is not fully offline.
  \item \todo{add the limitations the build actually reveals}
\end{itemize}

\section{Future Scope}

\begin{itemize}[nosep]
  \item More orders, products, and work centres, to test how negotiation
        behaves as contention grows.
  \item Richer negotiation, such as multi-round counter-offers or partial
        offer acceptance.
  \item Learned scoring weights derived from recorded outcomes, with the
        scoring rule kept inspectable.
  \item Integration with a real shop-floor data source instead of an injected
        disruption.
  \item \todo{add directions that emerge from the results}
\end{itemize}
